\chapter{Introduction}

Termination is a fundamental property of computer programs and a cornerstone of
formal verification. Infinite loops can lead to system failure, resource
exhaustion, and application crashes.
However, ensuring that a program will eventually halt, known as the Halting
Problem, was proved to be undecidable in the general case by Alan Turing
\cite{Turing1937}. Since no algorithm can universally solve the Halting Problem,
we must rely on restricted techniques that are sound however not complete.

When working with total functions, termination is a requirement for logical
consistency in formal verification and proof assistants.
Indeed, in such settings, non-terminating functions can be exploited to
effectively allow the type system to accept proofs of false statements
\cite{InfFalse}.
To illustrate this vulnerability in a system utilizing refinement types such as
\cite{RefType}, consider how an infinite loop can trick the compiler into
accepting a logical contradiction by forging an impossible type:
\begin{minted}{scala}
        def asNonEmpty[A](l: List[A]): {b: Boolean with b == true && l != Nil} =
          val res = asNonEmpty(l)
          res
        @main def main =
          val l = Nil: {l: List[Int] with asNonEmpty(l)}
\end{minted}
The non-terminating function \texttt{asNonEmpty} is considered well-typed,
and therefore the type checker accepts the impossible refinement type for
\texttt{Nil}. While the type system's logical soundness is compromised at
compile time, the impossible type is not materialized at runtime.

It is worth noting that there are consistent formal systems that do not require
termination to derive correct conclusions. They are usually based on partial
functions. Notable examples include the domain theory as mechanized in the LCF
proof assistant family \cite{LCFtoHOL}, and more recently, the Iris framework
\cite{Iris}. However, we want to work with total functions over algebraic
data types, as such termination remains a concern.

Consequently, proof assistants dealing with total
functions such as Rocq, Agda, and Lean employ techniques such as recursive calls
to be on syntactically smaller arguments, size-change principle
\cite{SizeChange}, well-founded measures \cite{wellfoundedRelation}, or manually
provided proofs.
For example, in a proof assistant like Rocq, termination is established when
recurring on structurally smaller components of an inductively defined type:
\begin{minted}{coq}
      Fixpoint sum (l: list nat): nat :=
        match l with
        | [] => 0
        | hd :: tl => hd + sum tl
        end.
\end{minted}

However, proof assistants are often too rigid for general-purpose programming.
Tools like Stainless \cite{DBLP:journals/corr/abs-1904-03482} provide a more
flexible approach by leveraging powerful SMT (Satisfiability Modulo Theories)
solvers, capable of proving complex termination properties.
Stainless excels at verifying semantic termination bounds over bounded and
unbound integers. For instance, Stainless can prove the termination of
\begin{minted}{scala}
        def sum(l: List[BigInt]): BigInt =
          if l.isEmpty then 0
          else l.head + sum(l.tail)
\end{minted}
But Stainless relies on external tools, and operates as an external toolchain to
Scala. This leaves a gap in general-purpose languages where developers often
lack a built-in, lightweight mechanism to verify the termination of critical
functions without relying on a separate verification framework.

To bridge this gap, this report presents an opt-in termination checker
implemented natively as a \texttt{MiniPhase} within the Scala 3 compiler.
It operates in a middle ground between the heavyweight, strict checks of proof
assistants and the lack of it in general-purpose languages.
Similar to proof assistants, it ensures termination natively through structural
decrease analysis over Algebraic Data Types:
\newpage
\begin{minted}{scala}
        @terminates
        def sum(list: List[Int]): Int =
          list match
            case Nil => 0
            case hd :: tl => sum(tl)
\end{minted}
The implementation focuses on common patterns Scala developers encounter,
however due to the lightweight structural analysis, it inevitably flags some
complex terminating functions as possibly non-terminating. To maintain usability, we provide several
escape hatches. Developers can use annotations like \texttt{@assumeTerminates},
\texttt{@assumeDecreases}, and \texttt{@assumeWellFounded} to manually assert
properties that the checker cannot automatically infer. This approach provides a
native and scalable way to enforce termination, falling back on developers
assertions when the automatic analysis limits are reached.

